\section{Algoritma RSA}

\begin{subcpmk}
  \item Mahasiswa mampu menganalisis derivasi matematis algoritma RSA, membuktikan koreksi dekripsinya melalui Teorema Euler, melakukan proses pembuatan kunci, enkripsi, dan dekripsi secara manual, serta mengevaluasi kompleksitas komputasinya.
\end{subcpmk}

\noindent\textit{Rujukan:} Munir \cite{munir2019,munir_slides_itb}; HAC Bab~8 \cite{hac_online}; Katz \& Lindell \cite{katz2020}; Boneh \cite{boneh_cs255}.

\subsection{Landasan Matematis: Fungsi Totient dan Teorema Euler}

RSA (Rivest-Shamir-Adleman) adalah algoritma kunci publik pertama yang mengimplementasikan konsep \textit{trapdoor one-way function} menggunakan operasi aritmetika modular. Keamanan RSA bersandar pada \textbf{Integer Factorization Problem} (IFP): fakta bahwa mengalikan dua bilangan prima besar sangat mudah, tetapi memfaktorkan hasil kalinya kembali menjadi dua prima asalnya secara komputasi mustahil untuk bilangan yang sangat besar.

Dua konsep matematika kunci yang mendasari RSA adalah:

\begin{enumerate}
    \item \textbf{Fungsi Totient Euler $\varphi(n)$}: Fungsi ini menghitung jumlah bilangan bulat positif kurang dari $n$ yang koprima dengan $n$. Jika $n = p \cdot q$ dan $p, q$ adalah bilangan prima, maka:
    \begin{equation}
      \varphi(n) = (p-1)(q-1)
    \end{equation}

    \item \textbf{Teorema Euler}: Jika $\gcd(a, n) = 1$, maka:
    \begin{equation}
      a^{\varphi(n)} \equiv 1 \pmod{n}
    \end{equation}
\end{enumerate}

\subsection{Prosedur Algoritma RSA}

Proses RSA terbagi menjadi tiga tahap: pembangkitan kunci, enkripsi, dan dekripsi.

\subsubsection{Pembangkitan Kunci (Key Generation)}
\begin{enumerate}
    \item Pilih dua bilangan prima besar $p$ dan $q$ secara acak dan independen.
    \item Hitung $n = p \cdot q$. Nilai $n$ disebut sebagai \textit{modulus} dan panjangnya dalam bit menentukan kekuatan kunci.
    \item Hitung $\varphi(n) = (p-1)(q-1)$.
    \item Pilih bilangan bulat $e$ sedemikian sehingga $1 < e < \varphi(n)$ dan $\gcd(e, \varphi(n)) = 1$. Nilai $e$ adalah \textit{public exponent}. Nilai standar yang umum digunakan adalah $e = 65537$ ($2^{16}+1$) karena efisien untuk komputasi.
    \item Hitung $d$ sebagai invers multiplikatif modular dari $e$ modulo $\varphi(n)$, sehingga:
    \begin{equation}
      e \cdot d \equiv 1 \pmod{\varphi(n)}
    \end{equation}
    Nilai $d$ dihitung menggunakan \textbf{Algoritma Euclidean Diperluas} (EEA).
\end{enumerate}
\textbf{Kunci Publik}: $(e, n)$ \\
\textbf{Kunci Privat}: $(d, n)$ atau cukup $(d)$

\subsubsection{Enkripsi}
Untuk pesan $M$ yang dinyatakan sebagai bilangan bulat $0 \le M < n$, cipherteks $C$ dihitung dengan:
\begin{equation}
  C = M^e \pmod{n}
\end{equation}

\subsubsection{Dekripsi}
Penerima menggunakan kunci privat $d$ untuk memulihkan pesan $M$ dari cipherteks $C$:
\begin{equation}
  M = C^d \pmod{n}
\end{equation}

\subsection{Pembuktian Koreksi RSA}

Untuk membuktikan bahwa dekripsi selalu mengembalikan pesan asli, kita harus menunjukkan bahwa $C^d \equiv M \pmod{n}$.

\begin{proof}
Berdasarkan definisi enkripsi:
\[ C^d \equiv (M^e)^d \equiv M^{ed} \pmod{n} \]
Karena $e \cdot d \equiv 1 \pmod{\varphi(n)}$, maka terdapat bilangan bulat $k$ sedemikian sehingga $ed = k\varphi(n) + 1$.
\[ M^{ed} = M^{k\varphi(n) + 1} = (M^{\varphi(n)})^k \cdot M^1 \]
Menurut Teorema Euler, $M^{\varphi(n)} \equiv 1 \pmod{n}$ (asumsi $M$ koprima dengan $n$):
\[ M^{ed} \equiv (1)^k \cdot M \equiv M \pmod{n} \]
Jika $M$ tidak koprima dengan $n$ (kasus yang sangat jarang untuk $p, q$ besar), pembuktian tetap berlaku menggunakan Teorema Sisa Tionghoa (CRT). Jadi, proses dekripsi selalu benar.
\end{proof}

\subsection{Contoh Perhitungan Manual (Small-Scale RSA)}

Misalkan kita ingin mengimplementasikan RSA dengan bilangan prima kecil:
\begin{enumerate}
    \item \textbf{Kunci}: Pilih $p = 11, q = 13 \rightarrow n = 143$.
    \item \textbf{Totient}: $\varphi(143) = (10)(12) = 120$.
    \item \textbf{Eksponen Publik}: Pilih $e = 7$ (cek $\gcd(7, 120) = 1 \checkmark$).
    \item \textbf{Eksponen Privat}: Cari $d$ sehingga $7d \equiv 1 \pmod{120}$.
    Menggunakan EEA: $120 = 17 \cdot 7 + 1 \rightarrow 1 = 120(1) - 17(7)$.
    Maka $d \equiv -17 \equiv 103 \pmod{120}$.
\end{enumerate}
Kunci Publik: $(7, 143)$, Kunci Privat: $(103, 143)$.

\textbf{Proses Enkripsi}: Pesan $M = 9$.
\[ C = 9^7 \pmod{143} = 4,782,969 \pmod{143} = 48 \]

\textbf{Proses Dekripsi}: Cipherteks $C = 48$.
\[ M = 48^{103} \pmod{143} = 9 \]

\subsection{Kompleksitas Komputasi dan Efisiensi}

Operasi utama dalam RSA adalah \textit{modular exponentiation} ($M^e \bmod n$). Jika dilakukan secara naif, kompleksitasnya adalah $\mathcal{O}(e)$, yang tidak praktis untuk $e$ besar. Oleh karena itu, digunakan algoritma \textbf{Square-and-Multiply} (atau \textit{Binary Exponentiation}) yang mereduksi kompleksitas menjadi $\mathcal{O}(\log e)$.

\begin{figure}[htbp]
  \centering
  \includegraphics[width=0.7\textwidth]{bab-06/rsa-square-multiply-diagram}
  \caption{Alur kerja algoritma Square-and-Multiply untuk eksponensiasi modular efisien.}
  \label{fig:rsa-square-multiply}
\end{figure}

\subsection{Keamanan dan Padding OAEP}

\textit{Textbook RSA} (RSA tanpa padding) memiliki kelemahan fatal:
\begin{enumerate}
    \item \textbf{Deterministik}: Pesan yang sama selalu menghasilkan cipherteks yang sama.
    \item \textbf{Kerentanan Matematis}: Jika $M$ kecil dan $e=3$, maka $M^3$ mungkin kurang dari $n$, sehingga $M$ dapat ditemukan dengan akar pangkat tiga biasa.
\end{enumerate}
Untuk mengatasinya, standar industri menggunakan \textbf{Optimal Asymmetric Encryption Padding (OAEP)}. OAEP menambahkan struktur acak (menggunakan fungsi hash) ke dalam pesan sebelum dienkripsi, sehingga mengubah RSA menjadi skema enkripsi probabilistik: pesan yang sama akan menghasilkan cipherteks berbeda setiap kali dienkripsi.

\begin{aktivitas}
  \item \textbf{Eksperimen RSA}: Implementasikan algoritma Square-and-Multiply dalam bahasa Python. Bandingkan waktu eksekusi antara metode perkalian berulang naif dengan metode Square-and-Multiply untuk eksponen $e > 10^6$.
\end{aktivitas}

\begin{latihan}
  \item Diketahui $p=17, q=19, e=5$. Hitunglah $n, \varphi(n), d$, kunci publik, dan kunci privat!
  \item Lakukan enkripsi dan dekripsi untuk pesan $M=10$ menggunakan parameter pada soal nomor 1!
  \item Mengapa pemilihan $e = 65537$ sangat populer dalam implementasi RSA? Jelaskan dari sisi efisiensi biner!
  \item Jika seorang penyerang berhasil menemukan $\varphi(n)$, apakah ia dapat memecahkan RSA? Jelaskan langkah-langkah yang akan dilakukan penyerang tersebut!
\end{latihan}
